\scanpage{41}
\mspage{2}

prenant pour $A$ l\deleted{\unsure{e}}a ``champ'' sur $X$ formé des $F$-torseurs à droite, la (Gr)-catégorie des aut\deleted{\unsure{b}}o-équivalences de $A$ avec lui-même s'interprète comme la catégorie des ``b\deleted{\unsure{t}}itorseurs'' sous $F$, i.e. des faisceaux \inserted{\handwritten{$P$}} sur lesquels $F$ opère à la fois à droite et à gauche, ces opérations commutant et chacune d'elles faisant de $P$ un torseur (à droite ou à gauche) sous $F$, - la composition $\otimes$ étant la composition de Baer évidente. Lorsque $F$ est encore de la forme $\underline O_X^{\handwritten{*}}$ ($\underline O_X$ étant un faisceau d'anneaux, qu'on ne suppose plus nécessairement commutatif) ces bitorseurs s'interprètent aussi en termes de bi-Modules\handwritten{``}inversibles\handwritten{''} sous $\underline O_X$.

\uline{Structure}.\deleted{\unsure{Les}} Soit $C$ une (Gr)-catégorie, on lui associe

a) le groupe $\handwritten{\pi}_0(C)$ des classes d'isomorphisme d'objets de $C$,

b) le groupe $\handwritten{\pi}_1(C)$ des automorphismes de $1_C$ (objet unité de $C$)

c) une action de $\handwritten{\pi}_0(C)$ sur $\handwritten{\pi}_1(C)$, en associant à tout objet $X$ de $C$ l'automorphisme \deleted{\unsure{de 1}} $\handwritten{\rho}(X)$ de $1_C$ déduit des deux isomorphismes
\[
\operatorname{Aut}(1_C)\handwritten{\rightrightarrows}\operatorname{Aut}(X)
\]
donnés par $u\handwritten{\mapsto}u\otimes\mathrm{id}_X$ et $u\handwritten{\longmapsto}\mathrm{id}_X\otimes u$. On prouve que $\handwritten{\pi}_1(C)$ est un groupe \uline{commutatif} et que l'on obtient bien par c) une structure de $\handwritten{\pi}_0(C)$-module sur celui-ci. Ceci posé, si on choisit pour tout $a\handwritten{\in\pi}_0(C)$ un représentant $L_a\handwritten{\in}\operatorname{Ob}C$ de $a$, et pour deux \deleted{objets} $a,b$ un isomorphisme
\[
\handwritten{\phi}_{a,b}: L_a\otimes L_b\handwritten{\simeq}L_{ab},
\]
alors pour trois éléments $a,b,c\handwritten{\in\pi}_0(C)$, l'isomorphisme d'associativité
\[
(L_a\otimes L_b)\otimes L_c\handwritten{\simeq}L_a\otimes(L_b\otimes L_c),
\]
compte tenu des isomorphismes $\handwritten{\phi}_{a,b}$, $\handwritten{\phi}_{ab,c}$, $\handwritten{\phi}_{b,c}$ et $\handwritten{\phi}_{a,bc}$, peut s'interpréter comme un isomorphisme
\[
L_{abc}\handwritten{\simeq}L_{abc},
\]
ou encore comme la tensorisation \deleted{\unsure{à}}\handwritten{à} gauche avec un élément bien déterminé
\[
f(a,b,c)\handwritten{\in\pi}_1(C)\deleted{\unsure{\times1_{\otimes}\times1_{\otimes}|\,1_C\otimes1_C}}.
\]
Les $\handwritten{\phi}_{a,b}$ étant choisis, on voit donc que la donnée d'un isomorphisme d'associativité fonctoriel $(L\otimes L')\otimes L''\handwritten{\simeq}L\otimes(L'\otimes L'')$ équivaut à la donnée de l'application
\[
f:\handwritten{\pi}_0\times\handwritten{\pi}_0\times\handwritten{\pi}_0\handwritten{\longrightarrow\pi}_1.
\]
On vérifie alors que l'axiome \deleted{\handwritten{\unsure{du pen}}}standard d'autocompatibilité d'une donnée d'associativité\inserted{\handwritten{(axiome du pentagone)}} s'exprime précisément par la condition que $f$ soit un $\handwritten{3}$-\uline{cocycle} du groupe $\handwritten{\pi}_{\deleted{\unsure{1}}0}$ à valeurs dans le groupe $\handwritten{\pi}_1$. D'autre part, l'indétermination dans le choix du système d'isomorp\deleted{\unsure{p}}hismes $\handwritten{\phi}_{a,b}$ est précisément donnée par une 2-cochaine arbitraire, et on voit que si on change $\handwritten{\phi}$ par une 2-cochaine $g$, $f$ est changé en $f+dg$ - donc l'ensemble des $f$ \deleted{\unsure{p}}co\deleted{\unsure{s}}rrespondants à des choix différents de $\handwritten{\phi}$ est exactement une classe de $\handwritten{3}$-cohomologie
% archive sheet number: 40
